\section{Implementation Details}
\label{app:details}

\subsection{Training and inference configuration}

\begin{table}[h]
\caption{Configuration used for the reported \method{} runs.}
\label{tab:implementation}
\centering
\small
\setlength{\tabcolsep}{6pt}
\begin{tabular}{ll}
\toprule
Setting & Value \\
\midrule
Backbone & Llama-3-8B-Instruct \\
Decomposer data & 12,000 examples; 4,000 per source \\
LoRA & rank 16, scale 32, dropout 0.05 \\
Target modules & all linear modules \\
Training & 3 epochs, bf16, effective batch 16 \\
Optimization & \(10^{-4}\), cosine decay, 10\% warmup \\
Maximum context & 1,024 tokens \\
Decoding & greedy; at most 48 answer tokens \\
Retrieval & BM25, top 8, Wikipedia-21M \\
Tail set & top \(k=5\) tokens \\
Step gate & \(\tau=0.002305\) \\
Direct verifier & \(\tau_d=0.0005\) \\
Post-decomposition verifier & \(\tau_p=0.0158\) \\
Type-aware margin & 2.5 \\
\bottomrule
\end{tabular}
\end{table}

\subsection{Prompts}

The answer probe uses the same answer-only instruction for direct and
sub-question generation:
\begin{quote}
\small
\emph{Answer the question with the answer phrase only. No explanation, no
preamble, no full sentence---just the answer.}
\end{quote}
For a decomposed step, the original question is prepended when it differs from
the resolved sub-question. Retrieved context, when available, appears between
the original question and sub-question.

The decomposer instruction asks for an ordered list of atomic sub-questions,
one per line, with the format \(\texttt{\#<index>:}\). Dependencies must use a
\(\texttt{\#N}\) placeholder. The final synthesizer receives the original
question and all sub-question/answer pairs. It returns only a short answer span
and contains explicit rules for comparison, yes/no, and year questions. The
supplementary source package contains the exact strings and chat-template
construction.

\subsection{Final routing rule}

The final rule is deterministic. It first handles exact candidate agreement,
empty or malformed outputs, unresolved placeholders, and obvious yes/no or
comparison-type mismatches. If these rules do not decide the route, it chooses
direct only when
\[
 \tl(a^{\mathrm{dir}})<\tau_d
 \quad\text{and}\quad
 \max_i\tl(a_i^{\mathrm{post}})>\tau_p
\]
or the decomposition is structurally broken; otherwise it chooses
decomposition. A final type-aware check is applied only to undecided
decomposition defaults. It scores answer-type compatibility, whether the final
answer repeats an intermediate answer, whether a step answer appears as an
entity fragment inside a later query, and whether a candidate simply repeats
part of the question. It switches to direct only when the direct score exceeds
the decomposition score by at least 2.5.

\section{Additional Retrieval Results}
\label{app:additional-results}

\subsection{Gold and untuned decompositions on MuSiQue}

The first 1,000 \musique{} development questions are all two-hop, allowing
comparison with the official decomposition. Table~\ref{tab:gold-decomp}
shows that \method{} improves over the original query but remains below a gold
chain. Replacing the bridge entity with the gold answer closes most of the gap,
indicating that bridge-answer drift is more important than sub-question wording
in this slice.

\begin{table}[h]
\caption{Answer-containing Hit@5 on the MuSiQue slice.}
\label{tab:gold-decomp}
\centering
\small
\setlength{\tabcolsep}{7pt}
\begin{tabular}{lc}
\toprule
Query condition & Hit@5 \\
\midrule
Original question & 19.0 \\
\method{} decomposition & 38.8 \\
Gold decomposition + gold bridge & \textbf{55.8} \\
Gold decomposition, unresolved \(\texttt{\#1}\) & 19.8 \\
\method{} wording + gold bridge & 54.1 \\
Untuned base wording + gold bridge & 52.9 \\
\bottomrule
\end{tabular}
\end{table}

\subsection{Title matching}

Original questions are long and constraint-rich, while decomposed queries are
short and entity-centered. Table~\ref{tab:title-policy} compares title matching
for both query types. These are live end-to-end reruns.

\begin{table}[h]
\caption{EM/F1 under three BM25 title policies.}
\label{tab:title-policy}
\centering
\small
\setlength{\tabcolsep}{7pt}
\begin{tabular}{lccc}
\toprule
Dataset & Titles off & Titles on & Path-aware \\
\midrule
\popqa{} & 51.9/54.0 & 52.4/54.6 & \textbf{53.4/55.7} \\
\hotpot{} & 37.3/48.7 & 37.1/48.4 & \textbf{37.9/49.3} \\
\musique{} & 23.3/31.6 & 24.9/\textbf{34.5} & \textbf{26.2}/34.4 \\
\twowiki{} & 44.5/50.4 & 49.1/55.5 & \textbf{49.6/56.3} \\
\bottomrule
\end{tabular}
\end{table}

\section{Additional Confidence Results}
\label{app:confidence}

\subsection{Signal definitions}

All scores are oriented so larger values indicate greater uncertainty.
\begin{itemize}
    \item \textbf{Maximum token entropy}: the largest next-token entropy in
    the generated answer.
    \item \textbf{Sequence NLL}: the summed negative log-likelihood of the
    generated tokens.
    \item \textbf{TARG margin}: the transformed average gap between the two
    largest logits over the first eight answer tokens.
    \item \textbf{Sample disagreement}: the fraction of ten sampled answers
    that differ from the greedy answer.
    \item \textbf{P(True)}: one minus the model's estimated probability that
    its proposed answer is true.
    \item \textbf{\tailleak{}}: the sum of negative log top-five mass in
    Equation~\ref{eq:tail-leakage}.
\end{itemize}
The first three and \tailleak{} use the same greedy pass. Sample disagreement
and P(True) require additional generations.

\subsection{Entropy components}

Table~\ref{tab:entropy-components} evaluates the three terms of
Equation~\ref{eq:entropy-decomposition}. The boundary-crossing component
carries most of the useful signal; the arrangement within the tail is close to
chance. These component values come from the saved diagnostic computation and
should be recomputed from complete logits before being elevated to a primary
claim.

\begin{table}[h]
\caption{Macro AUROC of entropy components.}
\label{tab:entropy-components}
\centering
\small
\begin{tabular}{lc}
\toprule
Component & Macro AUROC \\
\midrule
Leading set versus tail & 0.742 \\
Competition within top five & 0.712 \\
Arrangement within tail & 0.502 \\
Full token entropy & 0.726 \\
Answer-level \tailleak{} & \textbf{0.749} \\
\bottomrule
\end{tabular}
\end{table}

\subsection{End-to-end signal replacement}

\begin{table}[h]
\caption{EM/F1 with entropy and Tail Leakage.}
\label{tab:entropy-vs-tail}
\centering
\scriptsize
\setlength{\tabcolsep}{3.5pt}
\begin{tabular}{lccccc}
\toprule
Signal & \hotpot{} & \twowiki{} & \musique{} & \popqa{} & Macro \\
\midrule
Max entropy & 35.50/46.11 & \textbf{50.00}/\textbf{56.58} & 24.70/33.06 & 51.90/54.71 & 40.52/47.61 \\
\tailleak{} & \textbf{37.80}/\textbf{49.03} & 49.30/56.22 & \textbf{26.50}/\textbf{34.63} & \textbf{53.90}/\textbf{55.73} & \textbf{41.88}/\textbf{48.90} \\
\bottomrule
\end{tabular}
\end{table}

\subsection{Selective answering}

Table~\ref{tab:selective} ranks direct answers by \tailleak{} and retains only
the most confident fraction. Accuracy increases as coverage decreases, but the
hard multi-hop datasets remain unreliable even at low coverage.

\begin{table}[h]
\caption{Direct-path EM by confidence coverage.}
\label{tab:selective}
\centering
\small
\setlength{\tabcolsep}{6pt}
\begin{tabular}{lccccc}
\toprule
Coverage & \hotpot{} & \twowiki{} & \musique{} & \popqa{} & Pooled \\
\midrule
25\% & 57.6 & 34.0 & 22.4 & 79.2 & 52.8 \\
50\% & 45.4 & 26.8 & 15.2 & 57.0 & 38.9 \\
75\% & 38.0 & 22.1 & 12.5 & 46.3 & 30.2 \\
100\% & 30.3 & 18.0 & 10.2 & 37.8 & 24.1 \\
\bottomrule
\end{tabular}
\end{table}

\section{Routing and Failure Analysis}
\label{app:routing}

\subsection{Candidate overlap}

\begin{table}[h]
\caption{Correctness overlap before final routing.}
\label{tab:path-overlap}
\centering
\small
\setlength{\tabcolsep}{5pt}
\begin{tabular}{lrrrr}
\toprule
Dataset & Both correct & Direct only & Decomp. only & Neither \\
\midrule
\hotpot{} & 225 & 78 & 135 & 562 \\
\twowiki{} & 144 & 36 & 353 & 467 \\
\musique{} & 75 & 27 & 192 & 706 \\
\popqa{} & 346 & 32 & 194 & 428 \\
\bottomrule
\end{tabular}
\end{table}

\subsection{Verifier effect}

\begin{table}[h]
\caption{Verifier changes relative to always choosing decomposition.}
\label{tab:verifier-effect}
\centering
\small
\setlength{\tabcolsep}{5pt}
\begin{tabular}{lrrr}
\toprule
Dataset & Direct-only recovered & Decomp.-only lost & Net EM \\
\midrule
\hotpot{} & 39/78 & 21/135 & +1.80 \\
\twowiki{} & 3/36 & 7/353 & -0.40 \\
\musique{} & 9/27 & 11/192 & -0.20 \\
\popqa{} & 1/32 & 2/194 & -0.10 \\
\midrule
Pooled / macro & 52/173 & 41/874 & +0.28 \\
\bottomrule
\end{tabular}
\end{table}

\subsection{Decomposition-path failures}

Each incorrect multi-hop decomposition candidate is assigned one primary
cause. Evidence is counted as present when the required Wikipedia title
appears in the top ten results from either the original or resolved
sub-question queries.

\begin{table}[h]
\caption{Primary causes of incorrect decomposition answers.}
\label{tab:failure-taxonomy}
\centering
\scriptsize
\setlength{\tabcolsep}{3.5pt}
\begin{tabular}{lccc}
\toprule
Cause & \hotpot{} (\(N=640\)) & \twowiki{} (\(N=503\)) & \musique{} (\(N=733\)) \\
\midrule
Missing required evidence & 307 (48.0\%) & 311 (61.8\%) & 513 (70.0\%) \\
Wrong step answer & 131 (20.5\%) & 98 (19.5\%) & 73 (10.0\%) \\
Wrong final synthesis & 164 (25.6\%) & 92 (18.3\%) & 138 (18.8\%) \\
Lost original constraint & 38 (5.9\%) & 2 (0.4\%) & 9 (1.2\%) \\
\bottomrule
\end{tabular}
\end{table}

Missing evidence is the largest category on all three multi-hop datasets.
Routing cannot repair many of these cases because the direct candidate is also
wrong.

\subsection{Confidently wrong cases}

\begin{table}[h]
\caption{Tail Leakage when both candidates agree in correctness.}
\label{tab:quadrant-confidence}
\centering
\small
\setlength{\tabcolsep}{4.5pt}
\begin{tabular}{lcccc}
\toprule
Dataset & Both-correct mean & Both-wrong mean & AUROC & Confidently wrong \\
\midrule
\hotpot{} & 0.01260 & 0.12927 & 0.777 & 15.5\% \\
\twowiki{} & 0.03414 & 0.20215 & 0.751 & 20.1\% \\
\musique{} & 0.04945 & 0.25025 & 0.766 & 12.7\% \\
\popqa{} & 0.00974 & 0.10282 & 0.846 & 4.9\% \\
Pooled & 0.01877 & 0.17926 & 0.817 & 9.3\% \\
\bottomrule
\end{tabular}
\end{table}

Both-wrong examples have higher scores on average, but the groups overlap.
This is why \tailleak{} is used as a population-level ranking and gating signal
rather than a per-question correctness test.
